\begin{proof}
We first establish the two facts about Archimedean seminorms that will be
needed in the continuity argument.  Let \(R\) be a commutative ring with a
multiplicative seminorm \(x\), normalized by \(|1|_x=1\), and set
\(\mathfrak p=\ker(x)\) and
\(K=\operatorname{Frac}(R/\mathfrak p)\).  If \(|2|_x\leq 1\), write
\(n^m\) in base \(2\).  The number of nonzero binary digits is \(O(m)\),
whereas every power of \(2\) has absolute value at most \(1\).  The ordinary
triangle inequality gives \(|n|_x^m=O(m)\), and taking \(m\)-th roots shows
that \(|n|_x\leq 1\) for every positive integer \(n\).  Expanding
\((a+b)^m\), using
\(\bigl|\binom mj\bigr|_x\leq 1\), and then taking \(m\)-th roots gives
\[
 |a+b|_x\leq\max\{|a|_x,|b|_x\}.
\]
Thus \(x\) induces a non-Archimedean absolute value on \(K\).  The converse
follows from \(|2|_x=|1+1|_x\leq 1\).  In particular,
\(\rho(x)=0\) exactly on the non-Archimedean locus.

Suppose that \(|2|_x>1\).  The base-expansion proof of Ostrowski's
normalization gives, for integers \(a,b\geq 2\),
\[
 |a|_x^m\leq C_b(1+m\log_ba)
 \max\{1,|b|_x^{m\log_ba+O(1)}\},
 \qquad
 C_b:=\max_{0\leq j<b}|j|_x.
\]
Taking \(m\)-th roots and letting \(m\to\infty\), and then interchanging
\(a\) and \(b\), yields
\[
 |n|_x=n^{\rho(x)}\qquad(n\geq 1).
\]
The triangle inequality gives \(|2|_x\leq 2\), so
\(0<\rho(x)\leq 1\).  Put \(v(a):=|a|_x^{1/\rho(x)}\) on \(K\).  This is
multiplicative and restricts to the usual absolute value on \(\mathbf Q\).
It also satisfies the ordinary triangle inequality.  Indeed, with
\(A_0=v(a)\) and \(B_0=v(b)\), the binomial expansion gives
\[
 |a+b|_x^m
 \leq\sum_{j=0}^m
 \bigl(\tbinom mj A_0^jB_0^{m-j}\bigr)^{\rho(x)}
 \leq(m+1)(A_0+B_0)^{m\rho(x)}.
\]
Taking the power \(1/(m\rho(x))\) and letting \(m\to\infty\) proves
\(v(a+b)\leq A_0+B_0\).

Let \(\widehat K\) be the completion for \(v\).  The closure of
\(\mathbf Q\) in \(\widehat K\) is \(\mathbf R\), and \(\widehat K\) is a
commutative real Banach division algebra.  We recall the short
Gelfand--Mazur argument in the form needed here.  In a complex Banach
division algebra the spectrum of every element is nonempty: otherwise,
after applying continuous linear functionals, its resolvent would be a
bounded entire function tending to \(0\) at infinity, contrary to
Liouville's theorem.  A spectral value then forces the element to be a
complex scalar.  Hence every complex Banach division algebra is
\(\mathbf C\).  If a commutative real Banach division algebra \(E\)
contains a square root of \(-1\), this complex statement gives
\(E=\mathbf C\).  Otherwise \(E[T]/(T^2+1)\), equipped with any equivalent
complete finite-module algebra norm, is a complex Banach division algebra
and hence is \(\mathbf C\).  The closed subfield \(E\subset\mathbf C\)
contains \(\mathbf R\) but not \(\mathrm i\), and therefore equals
\(\mathbf R\).  It follows that \(\widehat K\) is isometrically
\(\mathbf R\) or \(\mathbf C\).  Composing
\(R\to K\to\widehat K\hookrightarrow\mathbf C\) produces a character
\(\chi_x\) satisfying
\begin{equation}\label{eq:hybrid-philippon-character}
 |a|_x=|\chi_x(a)|_\infty^{\rho(x)}.
\end{equation}

This character is unique up to conjugation.  In fact, suppose that
\(\sigma,\tau\colon R\to\mathbf C\) satisfy
\(|\sigma(a)|=|\tau(a)|\) for every \(a\).  They have the same kernel and
induce a modulus-preserving field isomorphism
\[
 \operatorname{Frac}(\sigma(R))\longrightarrow
 \operatorname{Frac}(\tau(R)),
 \qquad
 \frac{\sigma(a)}{\sigma(b)}\longmapsto
 \frac{\tau(a)}{\tau(b)}.
\]
It extends isometrically to the closures \(L,M\subset\mathbf C\).  A closed
subfield of \(\mathbf C\) containing \(\mathbf Q\) contains \(\mathbf R\);
it equals \(\mathbf R\) if it has no nonreal element and equals
\(\mathbf C\) otherwise.  In the real case the extension is the identity.
In the complex case it is a continuous automorphism of \(\mathbf C\)
fixing \(\mathbf R\), so it sends \(\mathrm i\) to \(\mathrm i\) or
\(-\mathrm i\) and is the identity or complex conjugation.  This proves the
asserted uniqueness in \eqref{eq:hybrid-philippon-character}.

We next record the corresponding topology of characters.  Let \(B\) be a
finitely generated \(\mathbf Z\)-algebra, equipped with the pointwise
topology on its multiplicative seminorms, and write
\[
 B=\mathbf Z[t_1,\ldots,t_r]/I.
\]
Suppose that a net of Archimedean seminorms \(x_\lambda\to x\) satisfies
\(\rho(x_\lambda)\geq\varepsilon>0\), and choose arbitrary representing
characters \(\chi_\lambda\).  Since
\(|2|_{x_\lambda}\to|2|_x\), one has
\(\rho(x_\lambda)\to\rho(x)\).  The values
\(|t_j|_{x_\lambda}\) are bounded, say by \(K_j\), and hence
\[
 |\chi_\lambda(t_j)|
 =|t_j|_{x_\lambda}^{1/\rho(x_\lambda)}
 \leq\max\{1,K_j^{1/\varepsilon}\}.
\]
The points
\((\chi_\lambda(t_1),\ldots,\chi_\lambda(t_r))\) therefore lie in one
compact polydisc.  From every subnet we may extract a further subnet
converging to \(z\).  Every polynomial in \(I\) vanishes at \(z\), so
evaluation at \(z\) defines a character
\(\eta\colon B\to\mathbf C\).  If \(b\in B\) is represented by a polynomial
\(P\), then
\[
 \chi_\lambda(b)=P(\chi_\lambda(t_1),\ldots,
 \chi_\lambda(t_r))\longrightarrow P(z)=\eta(b).
\]
Consequently
\[
 |b|_x=\lim_\lambda|b|_{x_\lambda}
 =\lim_\lambda|\chi_\lambda(b)|^{\rho(x_\lambda)}
 =|\eta(b)|^{\rho(x)}.
\]
The uniqueness just proved shows that \(\eta\) is \(\chi_x\) or
\(\overline{\chi_x}\).  Thus every subnet of the chosen characters has a
further subnet converging pointwise on \(B\) to one of these two
characters.

We now turn to the Archimedean polynomial functional.  For a nonzero
multihomogeneous complex polynomial \(F\) in fixed blocks and fixed degrees,
D'Andrea--Krick--Sombra prove
\[
 \bigl|m(F)-\log\|F\|_{\max}\bigr|
 \leq\sum_i\log(N_i+1)d_i
\]
and
\[
 0\leq\operatorname{ph}(F)-m(F)
 \leq\sum_i\log(N_i+1)d_i;
\]
see \cite[Lemma~2.32, Definition~2.33, and (2.34)--(2.35)]{DKS13Heights}.
Here \(m(F)\) is the grouped-variable Mahler measure and
\(\|F\|_{\max}\) is the maximum coefficient modulus.  Thus, with the
harmless fixed-type constant
\begin{equation}\label{eq:hybrid-philippon-fixed-bound}
 C:=2\sum_i\log(N_i+1)d_i,
\end{equation}
one has
\begin{equation}\label{eq:hybrid-philippon-coefficient-bound}
 \bigl|\operatorname{ph}(F)-\log\|F\|_{\max}\bigr|\leq C.
\end{equation}

We also need continuity in the coefficients.  Zhang proves continuity and
positivity of the Chow metric, and hence of its logarithmic spherical
integral, on the nonzero projective coefficient space of a fixed equal
Chow type \cite[Section~1.5 and Theorem~3.6(1)]{Zhang96Semistable}.
The reduction from arbitrary fixed block sizes and multidegrees is
elementary.  If every \(d_i=0\), the coefficient space is one-dimensional
and continuity follows directly from
\(\operatorname{ph}(a)=\log|a|_\infty\).  Otherwise put
\(Q=\max_iN_i\) and \(D=\max_i d_i>0\).  Pad the \(i\)-th
block with \(Q-N_i\) unused variables and send
\begin{equation}\label{eq:hybrid-philippon-padding}
 F\longmapsto\iota(F):=F\prod_iU_{i0}^{D-d_i}.
\end{equation}
This is an injective linear map into the equal-type coefficient space.  On
the unit sphere in \(\mathbf C^Q\), write the first \(N_i\) coordinates as
\(r_i\xi_i\), where \(\xi_i\) is uniformly distributed on the unit sphere
in \(\mathbf C^{N_i}\) and is independent of \(r_i\).  Multihomogeneity
gives, off a null set,
\[
 \log|\iota(F)|=
 \log|F(\xi_1,\ldots,\xi_m)|
 +\sum_i\bigl(D\log r_i+(D-d_i)\log|\xi_{i0}|\bigr).
\]
The added logarithms are integrable and their integrals depend only on the
fixed type.  After including the fixed harmonic corrections, we obtain
\[
 \operatorname{ph}_{\mathrm{equal}}(\iota(F))
 =\operatorname{ph}(F)+\kappa
\]
with \(\kappa\) independent of \(F\).  Zhang's result therefore proves
projective coefficient continuity for the original type.  Homogeneity of
the normalized sphere integral lifts this to the punctured coefficient
vector space.  The same formula also gives
\begin{equation}\label{eq:hybrid-philippon-complex-homogeneity}
 \operatorname{ph}(\lambda F)=\log|\lambda|_\infty+
 \operatorname{ph}(F),
 \qquad
 \operatorname{ph}(\overline F)=\operatorname{ph}(F).
\end{equation}
In particular, \(\operatorname{ph}(F)\) is finite, continuous, homogeneous,
and invariant under complex conjugation on every punctured fixed-type
coefficient space.

Continuity is local on \(U\).  We may therefore work on a Banach analytic
domain \(V\subset U\) on which all coefficients belong to one Banach
algebra \(A_V\).  If a rational localization is used, include its finitely
many denominators and their inverses in what follows; they are units on
\(V\), so evaluation of the inverse is continuous.  Let \(B\subset A_V\)
be the finitely generated \(\mathbf Z\)-algebra generated by the
coefficients and these localization elements.  Define
\begin{equation}\label{eq:hybrid-philippon-max-function}
 M(x):=\max_{\alpha\in S}
 \bigl(\log|c_\alpha(x)|_x+w_\alpha\bigr)
 =\log\max_{\alpha\in S}
 \bigl(\mathrm e^{w_\alpha}|c_\alpha(x)|_x\bigr).
\end{equation}
The quantity inside the last logarithm is a finite maximum of continuous
nonnegative functions and is everywhere positive by hypothesis.  Hence
\(M\) is a finite continuous function, even where some coefficients vanish.

If \(\rho(x)>0\), the character construction applied to \(A_V\) gives
\(\chi_x\).  Only its restriction to \(B\) enters
\eqref{eq:hybrid-philippon-arch-polynomial}.  The polynomial \(F_x\) is
nonzero.  Changing \(\chi_x\) conjugates all coefficients simultaneously,
so \eqref{eq:hybrid-philippon-complex-homogeneity} shows that
\eqref{eq:hybrid-philippon-arch-value} is independent of this choice.

Let \(x_\lambda\to x\) with \(\rho(x)>0\).  Eventually
\(\rho(x_\lambda)\geq\rho(x)/2\), and all these points are Archimedean.
Choose arbitrary representing characters.  The finite-type character
argument shows that every subnet has a further subnet whose characters on
\(B\) converge to \(\chi_x\) or \(\overline{\chi_x}\).  Along it, for each
\(\alpha\),
\[
 \chi_{x_\lambda}(c_\alpha)
 \mathrm e^{w_\alpha/\rho(x_\lambda)}
 \longrightarrow
 \chi_x(c_\alpha)\mathrm e^{w_\alpha/\rho(x)}
\]
or all limits are conjugated.  Coefficient continuity and conjugation
invariance of \(\operatorname{ph}\) give
\[
 \rho(x_\lambda)\operatorname{ph}(F_{x_\lambda})
 \longrightarrow\rho(x)\operatorname{ph}(F_x).
\]
A net converges to a point when every subnet has a further subnet converging
to that point: otherwise one can choose a subnet outside a fixed
neighbourhood.  Thus \(\Phi\) is continuous on the locus \(\rho>0\).

It remains to cross its boundary.  For \(\rho(x)>0\), one has the exact
identity
\begin{align}
 \rho(x)\log\|F_x\|_{\max}
 &=\rho(x)\max_\alpha
 \left(\log|\chi_x(c_\alpha)|_\infty+
 \frac{w_\alpha}{\rho(x)}\right)\notag\\
 &=\max_\alpha\bigl(\log|c_\alpha(x)|_x+w_\alpha\bigr)
 =M(x).
 \label{eq:hybrid-philippon-exact-max}
\end{align}
This remains valid when some coefficients are zero, with
\(\log0=-\infty\).  Combining
\eqref{eq:hybrid-philippon-coefficient-bound} and
\eqref{eq:hybrid-philippon-exact-max} yields
\begin{equation}\label{eq:hybrid-philippon-boundary-estimate}
 |\Phi(x)-M(x)|\leq C\rho(x).
\end{equation}
If \(x_\lambda\to x\) and \(\rho(x)=0\), the left side of
\eqref{eq:hybrid-philippon-boundary-estimate} is zero at every
non-Archimedean \(x_\lambda\), while the displayed estimate applies to
every Archimedean \(x_\lambda\).  Since \(M(x_\lambda)\to M(x)\) and
\(\rho(x_\lambda)\to0\), we obtain
\(\Phi(x_\lambda)\to M(x)=\Phi(x)\).  Conversely, continuity of \(\rho\)
prevents a net with \(\rho=0\) from converging to a point with \(\rho>0\).
This proves continuity everywhere.

If \(\lambda\ne0\), then on the non-Archimedean locus
\[
 \Phi(\lambda F)=\max_\alpha
 \bigl(\log|\lambda c_\alpha|+w_\alpha\bigr)
 =\log|\lambda|+\Phi(F).
\]
On the Archimedean locus,
\eqref{eq:hybrid-philippon-complex-homogeneity} and
\eqref{eq:hybrid-philippon-character} give the same formula.  This proves
\eqref{eq:hybrid-philippon-homogeneity}.  Finally,
\(eF_e=e'F_{e'}=ueF_{e'}\), so \(F_{e'}=u^{-1}F_e\), and homogeneity gives
\eqref{eq:hybrid-philippon-frame-change} with the stated minus sign.

We finish with the basic boundary and failure tests.  On the hybrid spectrum
of \(\mathbf Z\), take the degree-one form \(U_0+U_1+U_2\) and zero weights.
Then \(M=0\), and \eqref{eq:hybrid-philippon-boundary-estimate} gives
\(|\Phi|\leq C\rho\), so the Archimedean values converge to the Gauss value
as \(\rho\downarrow0\).  A tuple such as \((t,1)\) may specialize to
\((0,1)\) without affecting continuity, whereas allowing all coefficients
to vanish would already make the value of \(tU_0\) tend to \(-\infty\).
For a real Archimedean character, the sphere integral is that of its
complexification and is unchanged by the two conjugate representatives.
No sign condition on \(w_\alpha\) entered either
\eqref{eq:hybrid-philippon-exact-max} or
\eqref{eq:hybrid-philippon-boundary-estimate}.  In degree zero with one
coefficient, the frame calculation is simply
\(\log|u^{-1}r|=\log|r|-\log|u|\), confirming the sign once more.
\end{proof}
